\documentclass[12pt]{article}

\usepackage[T2A]{fontenc}
\usepackage[utf8]{inputenc}
\usepackage[russian]{babel}
\usepackage{amsmath}
\usepackage{amssymb}
\usepackage{amsthm}

\theoremstyle{plain}
\newtheorem{theorem}{Теорема}
\newtheorem{lemma}{Лемма}
\newtheorem{proposition}{Предложение}
\newtheorem{corollary}{Следствие}

\theoremstyle{definition}
\newtheorem{definition}{Определение}

\theoremstyle{remark}
\newtheorem{remark}{Замечание}

\newcommand{\RR}{\mathbb{R}}
\newcommand{\ZZ}{\mathbb{Z}}
\DeclareMathOperator{\diam}{diam}
\DeclareMathOperator{\dist}{dist}

\title{Периодические раскраски евклидовых пространств по решётчатым
разбиениям Вороного:\\ метод и вычислительные результаты}
\author{}
\date{}

\begin{document}

\maketitle

\begin{abstract}
Описан метод построения периодических раскрасок пространства $\RR^n$,
запрещающих замкнутый интервал расстояний: пространство разбивается на ячейки
Вороного решётки $\Lambda$, а цвета назначаются по классам смежности
подрешётки $\Gamma\subset\Lambda$ индекса $k$. Метод опирается на две леммы ---
о центральной симметрии ячеек и лемму Иванова о середине,
$D(v)=2\dist(v/2,V_0)$, --- и даёт вычислимый сертификат: раскраска в $k$
цветов запрещает интервал $[1,\ell]$ при $\ell\le d=D/\diam V_0$. Приведены
точные значения констант $d$, проверенные двумя независимыми реализациями, в
размерностях $n=2,3,4,5,6,8$, а также закон
$d=\sqrt{7/3}\,/\,(2R/\lambda_1)$ для подрешёток вида $(3+\omega)\Lambda$
эйзенштейновых решёток.
\end{abstract}

\section{Введение}

Проблема Нельсона---Хадвигера состоит в определении хроматического числа
$\chi(\RR^n)$ --- наименьшего числа цветов, которыми можно раскрасить точки
евклидова пространства $\RR^n$ так, чтобы любые две точки на расстоянии $1$
получили разные цвета. Задача была поставлена в середине XX века и открыта до
сих пор даже для плоскости; её истории посвящены
\cite{Jensen1995,Soifer1991,Soifer2004,Soifer2008,Raigorodskii2003}, один из
первых родственных результатов --- теорема Хадвигера о покрытиях
\cite{Hadwiger1944}; задача тесно связана и с другими классическими проблемами
комбинаторной геометрии, например с проблемой Борсука \cite{Borsuk1933}.

В малых размерностях известно:
\[
5 \le \chi(\RR^2) \le 7, \qquad
6 \le \chi(\RR^3) \le 15, \qquad
9 \le \chi(\RR^4) \le 49.
\]
Нижние оценки принадлежат де Грею \cite{deGrey2018} ($n=2$), Нечуштану
\cite{Nechushtan2002} ($n=3$) и Эксу, Исмаилеску и Лиму \cite{Exoo2014}
($n=4$); верхние даются раскраской плоскости шестиугольниками (Исбелл, см.
\cite{Soifer2008}), раскраской Кулсона \cite{Coulson2002} ($n=3$) и
решётчатыми конструкциями Иванова \cite{Ivanov2006} и Армана, Бондаренко,
Прымака и Радченко \cite{Arman2021} ($n=4$); оценка $\chi(\RR^4)\le 54$ была
ранее получена Радоичичем и Тотом \cite{Radoicic2003}. Асимптотически
\[
(1.239+o(1))^n \le \chi(\RR^n) \le (3+o(1))^n
\]
(нижняя оценка --- \cite{Raigorodskii2000}, усиливающая метод Франкла---Уилсона
\cite{Frankl1981}; верхняя --- Ларман и Роджерс \cite{Larman1972}).

Мы рассматриваем интервальное обобщение задачи, систематически изученное
Ивановым \cite{Ivanov2011}: запрещается не одно расстояние, а целый отрезок
$[1,\ell]$, $\ell\ge 1$. Соответствующее хроматическое число обозначается
$\chi(\RR^n;[1,\ell])$. Так как $1\in[1,\ell]$, любая верхняя оценка
интервального хроматического числа является верхней оценкой и для
$\chi(\RR^n)$; все перечисленные выше верхние оценки малых размерностей
получаются именно интервальными решётчатыми конструкциями. Отметим и смежное,
но отличное по постановке направление: хроматические числа графов Вороного
решёток, где запрет накладывается на пары касающихся ячеек (комбинаторный
инвариант разбиения, а не метрики), изучены в \cite{DutourSikiric2021}.

В настоящей работе изложен метод построения раскрасок по разбиениям
Вороного решёток (методология Иванова \cite{Ivanov2011}), описан алгоритм с
сертификатом полноты перебора и приведены вычислительные результаты: точные
значения нормированных констант $d$ для проверенных конструкций в размерностях
$2$, $3$, $4$, $5$, $6$, $8$, включая константы конструкций из
\cite{Arman2021}, ранее в явном виде не публиковавшиеся.

\section{Определения и обозначения}

\begin{definition}
Пусть $(X,\rho)$ --- метрическое пространство, $A\subset\RR_{+}$ --- множество
запрещённых расстояний. \emph{Хроматическим числом} $\chi((X,\rho);A)$
называется наименьшее число цветов $k$, при котором существует разбиение
\[
X = V_1 \sqcup V_2 \sqcup \dots \sqcup V_k,
\]
такое, что $\rho(x,y)\notin A$ для любых $x,y\in V_i$, $1\le i\le k$. При
$A=\{1\}$ пишем просто $\chi(X)$, при $A=[1,\ell]$ --- $\chi(X;[1,\ell])$.
\end{definition}

\begin{remark}
Гомотетии переводят правильные раскраски в правильные, поэтому
$\chi(\RR^n;[a,b])$ зависит только от отношения $b/a$; интервал всегда можно
привести к виду $[1,\ell]$.
\end{remark}

\begin{definition}
\emph{Решёткой} $\Lambda\subset\RR^n$ называется дискретная подгруппа,
порождённая целочисленными комбинациями базисных векторов
$\mathbf{a}_1,\dots,\mathbf{a}_n$:
\[
\Lambda=\Bigl\{\sum_{i=1}^{n}k_i\mathbf{a}_i \;\Bigm|\; k_i\in\ZZ\Bigr\}.
\]
Определитель решётки $\det\Lambda=|\det(\mathbf{a}_1,\dots,\mathbf{a}_n)|$ ---
объём фундаментального параллелепипеда.
\end{definition}

\begin{definition}
\emph{Подрешётка} $\Gamma\subseteq\Lambda$ --- подгруппа, сама являющаяся
решёткой полного ранга. Её базис выражается через базис $\Lambda$
целочисленной невырожденной матрицей $M$ (\emph{матрицей перехода}; строки
$M$ --- координаты базисных векторов $\Gamma$ в базисе $\Lambda$).
\emph{Индекс} подрешётки
\[
[\Lambda:\Gamma] \;=\; |\det M| \;=\; \frac{\det\Gamma}{\det\Lambda}
\]
равен числу классов смежности факторгруппы $\Lambda/\Gamma$.
\end{definition}

\begin{definition}
\emph{Ячейкой Вороного} (областью Вороного) точки $\mathbf{a}\in\Lambda$
называется множество
\[
V(\mathbf{a})=\{\,\mathbf{x}\in\RR^n \mid
\|\mathbf{x}-\mathbf{a}\|\le\|\mathbf{x}-\mathbf{b}\|
\ \ \forall\,\mathbf{b}\in\Lambda\,\};
\]
центральную ячейку $V(\mathbf{0})$ обозначаем $V_0$. Ячейки покрывают
$\RR^n$ и пересекаются лишь по границам; $\diam V_0$ --- диаметр ячейки.
\end{definition}

\begin{remark}\label{rem:relevant}
$V_0$ --- пересечение полупространств
$\langle\mathbf{x},v\rangle\le\|v\|^2/2$ по $v\in\Lambda\setminus\{0\}$,
причём достаточно конечного набора \emph{релевантных} векторов (тех, что
задают гиперграни). По теореме Вороного вектор $v$ релевантен тогда и только
тогда, когда $\pm v$ --- единственные векторы минимальной нормы в классе
$v+2\Lambda$; поэтому релевантные векторы находятся сертифицированно ---
решением задачи ближайшего вектора в каждом из $2^n-1$ нетривиальных классов
$\Lambda/2\Lambda$, и их не более $2(2^n-1)$.
\end{remark}

\emph{Периодическая раскраска.} Пусть $\Gamma\subset\Lambda$ --- подрешётка
индекса $k$. Точки ячейки $V(\mathbf{a})$, $\mathbf{a}\in\Lambda$, красим
цветом класса смежности $\mathbf{a}+\Gamma$; получается $k$ цветов. Точки
одного цвета лежат либо в одной ячейке, либо в ячейках $V(\mathbf{a})$,
$V(\mathbf{b})$ с $\mathbf{a}-\mathbf{b}\in\Gamma$; поэтому ключевая величина
--- расстояние между ячейками, сдвинутыми на вектор подрешётки:
\[
D(v)=\dist\bigl(V_0,\,V(v)\bigr),\qquad
D=D(\Gamma)=\min_{v\in\Gamma\setminus\{0\}} D(v).
\]

\section{Две леммы}

\begin{lemma}[классические свойства ячеек, см.~\cite{Voronoi1952}]
\label{lem:sym}
Ячейка Вороного $V_0$ решётки $\Lambda\subset\RR^n$ --- выпуклый многогранник,
центрально-симметричный относительно начала координат: $V_0=-V_0$. Все ячейки
--- трансляты центральной: $V(\mathbf{a})=\mathbf{a}+V_0$ для
$\mathbf{a}\in\Lambda$. Гиперграни $V_0$ также центрально-симметричны
(каждая --- относительно своего центра).
\end{lemma}

\begin{lemma}[лемма Иванова о середине, \cite{Ivanov2011}]
\label{lem:midpoint}
Пусть $v\in\Lambda\setminus\{0\}$ и $s=v/2$. Тогда
\[
D(v)=\dist\bigl(V_0,\,V(v)\bigr)=2\,\dist(s,\,V_0),
\]
причём минимум расстояния между $V_0$ и $V(v)$ реализуется на паре точек,
симметричных относительно $s$.
\end{lemma}

\begin{proof}[Набросок доказательства]
Пусть $\sigma\colon x\mapsto v-x$ --- центральная симметрия относительно
середины $s$ отрезка между центрами ячеек. По лемме~\ref{lem:sym}
\[
\sigma(V_0)=v-V_0=v+V_0=V(v),
\]
т.е. $\sigma$ меняет местами $V_0$ и $V(v)$. Пусть пара $(p,q)$, $p\in V_0$,
$q\in V(v)$, реализует $D(v)$; тогда пара $(\sigma(q),\sigma(p))$ реализует его
же. По выпуклости
\[
p'=\tfrac12\bigl(p+\sigma(q)\bigr)\in V_0,\qquad
q'=\tfrac12\bigl(q+\sigma(p)\bigr)\in V(v),
\]
причём прямое вычисление даёт $p'-q'=p-q$ и $q'=\sigma(p')$. Итак, минимум
реализуется на паре, симметричной относительно $s$, а для такой пары
$\|p'-\sigma(p')\|=2\|p'-s\|\ge 2\dist(s,V_0)$. Обратно, если $m\in V_0$ ---
ближайшая к $s$ точка, то $\sigma(m)\in V(v)$ и
$\|m-\sigma(m)\|=2\|m-s\|=2\dist(s,V_0)$, откуда $D(v)\le 2\dist(s,V_0)$.
\end{proof}

Лемма~\ref{lem:midpoint} сводит вычисление расстояния между двумя
многогранниками к расстоянию от одной точки до одного многогранника ---
центральной ячейки $V_0$.

\section{Нормировка и корректность раскраски}

\begin{remark}\label{rem:diam}
Из центральной симметрии следует
$\diam V_0 = 2\max_{x\in V_0}\|x\| = 2R$, где $R$ --- радиус покрытия решётки:
действительно, $\|x-y\|\le\|x\|+\|y\|\le 2\max\|\cdot\|$, а для точки $x$
максимальной нормы пара $(x,-x)$ реализует равенство. В частности, все
диаметральные пары точек $V_0$ антиподальны.
\end{remark}

Нормированное запрещённое расстояние определяется как
\[
d \;=\; \frac{D}{\diam V_0}
\]
--- именно в таком виде, без каких-либо дополнительных множителей: двойка из
леммы~\ref{lem:midpoint} уже входит в $D=2\dist(v/2,V_0)$.

\begin{theorem}\label{thm:main}
Пусть $\Gamma\subset\Lambda$ --- подрешётка индекса $k$,
$D=\min_{v\in\Gamma\setminus\{0\}}D(v)$ и $d=D/\diam V_0\ge 1$. Тогда
\[
\chi\bigl(\RR^n;[1,\ell]\bigr)\;\le\;k
\qquad\text{для всех}\quad 1\le\ell\le d.
\]
\end{theorem}

\begin{proof}[Набросок доказательства]
Применим гомотетию с коэффициентом $1/\diam V_0$; после неё $\diam V_0=1$, а
расстояние между одноцветными ячейками не меньше $d$. Пусть $x,y$ --- точки
одного цвета. Если они в одной ячейке, то $\|x-y\|\le 1$; если в разных, то
$\|x-y\|\ge d\ge\ell$. Значит, монохромных пар с расстоянием строго внутри
интервала $(1,\ell)$ нет вовсе. Остаются граничные реализации: расстояние
ровно $1$ (диаметральные пары одной ячейки) и, при $\ell=d$, расстояние ровно
$\ell$ (ближайшие пары одноцветных ячеек). Обе реализуются только парами
граничных точек и устраняются аккуратным распределением границ между ячейками
(«полуоткрытые ячейки»); существенно, что по замечанию~\ref{rem:diam}
диаметральные пары антиподальны, а по лемме~\ref{lem:midpoint} ближайшие пары
симметричны относительно середины. Подробности см. у Кулсона
\cite{Coulson2002} и в \cite[Prop.~5]{Arman2021}.
\end{proof}

\begin{remark}[граничный случай $d=1$]
Конструкция содержательна тогда и только тогда, когда $d\ge 1$. При $d=1$
запрещённый интервал вырождается в точку $[1,1]$, а $D=\diam V_0$: оба
граничных явления из доказательства совпадают, и корректность раскраски
целиком держится на раскраске границ полуоткрытыми ячейками. Таков, например,
классический случай $15$ цветов в $\RR^3$ (решётка BCC, $D=\diam
V_0=\sqrt{5}$ в стандартных координатах) \cite{Coulson2002}.
\end{remark}

\section{Алгоритм}

\subsection{Перебор подрешёток: эрмитова нормальная форма}

Каждая подрешётка $\Gamma\subset\Lambda$ индекса $k$ имеет единственную
матрицу перехода $M$ в \emph{эрмитовой нормальной форме} (HNF): $M$
верхнетреугольна, её диагональ $d_1,\dots,d_n>0$, $\prod_j d_j=k$, а элементы
над диагональю в $j$-м столбце приведены по модулю $d_j$. Отсюда число
подрешёток индекса $k$ в решётке ранга $n$ равно
\[
N_n(k)\;=\;\sum_{\substack{d_1\cdots d_n=k\\ d_j\ge 1}}\;
\prod_{j=1}^{n} d_j^{\,j-1};
\qquad
N_4(k)=\sum_{d_1d_2d_3d_4=k} d_2\,d_3^2\,d_4^3,
\]
например $N_4(49)=140\,050$. Подчеркнём: суммирование идёт по \emph{всем
упорядоченным} наборам диагоналей $(d_1,\dots,d_n)$; сведение к неубывающим
наборам теряет подрешётки, в общем случае неэквивалентные учтённым.

\subsection{Полнота перебора кандидатов}

\begin{proposition}\label{prop:bound}
Для любого $v\in\Lambda\setminus\{0\}$ выполнено
$D(v)\ \ge\ \|v\|-\diam V_0$.
\end{proposition}

\begin{proof}
Для любого $m\in V_0$ имеем (замечание~\ref{rem:diam})
$\|s-m\|\ge\|s\|-\|m\|\ge\tfrac12\|v\|-\tfrac12\diam V_0$, где $s=v/2$;
остаётся применить лемму~\ref{lem:midpoint}.
\end{proof}

\begin{corollary}\label{cor:window}
Пусть $D_{\text{тек}}$ --- текущий рекорд минимума в переборе векторов
подрешётки $\Gamma$. Вектор $v$ с
$\|v\|\ \ge\ D_{\text{тек}}+\diam V_0$
не может уменьшить минимум. Следовательно, точное перечисление всех векторов
$\Gamma$ в шаре $\|v\|<D_{\text{тек}}+\diam V_0$ даёт сертификат полноты
вычисления $D(\Gamma)$.
\end{corollary}

\subsection{Описание алгоритма}

Для решётки $\Lambda$ (базис задаётся пользователем) и индекса $k$:

\begin{enumerate}
\item Привести базис $\Lambda$ алгоритмом LLL; построить $V_0$ по релевантным
      векторам (замечание~\ref{rem:relevant}: сертифицированный перебор по
      классам $\Lambda/2\Lambda$); выполнить самопроверку ячейки (объём равен
      $\det\Lambda$, центральная симметрия, соотношение Эйлера для
      $f$-вектора); вычислить $\diam V_0$.
\item Перечислить все $N_n(k)$ подрешёток индекса $k$ в HNF (все упорядоченные
      диагонали и все наддиагональные элементы).
\item Для каждой подрешётки $\Gamma$: LLL-привести её базис и точно
      перечислить кандидатов $v\in\Gamma$, $0<\|v\|<D_{\text{тек}}+\diam V_0$
      (следствие~\ref{cor:window}); граница сжимается по мере уменьшения
      рекорда.
\item Для каждого кандидата $v$ вычислить $D(v)=2\dist(v/2,V_0)$
      (лемма~\ref{lem:midpoint}); расстояние от точки до многогранника
      считается точным геометрическим примитивом (каскад проекций на грани
      либо алгоритм GJK по вершинам).
\item Положить $D(\Gamma)=\min_v D(v)$; по всем подрешёткам выбрать
      $D=\max_\Gamma D(\Gamma)$ (ветви отсекаются, как только промежуточный
      минимум опускается ниже текущего рекорда).
\item Выдать $d=D/\diam V_0$, матрицу перехода оптимальной подрешётки (HNF в
      базисе, переданном пользователем) и её LLL-приведённый базис в
      объемлющих координатах.
\end{enumerate}

Метод реализован дважды независимо: эталонным пакетом на python (ячейка ---
пересечение полупространств релевантных векторов, расстояние --- каскад
проекций на грани) и ядром на C++ (релевантные векторы ---
CVP-перебор по классам $\Lambda/2\Lambda$, расстояние --- GJK по вершинам).
Совпадение результатов двух реализаций с различными геометрическими
примитивами служит взаимной проверкой; флагманский результат ($D_4$, $k=49$)
дополнительно сверен с независимым QP-решателем (совпадение до девяти знаков,
$6000$ адверсариальных сверок расстояний без единого расхождения).

\section{Вычислительные результаты}

\subsection{Проверенные интервальные константы}

В таблице~\ref{tab:constants} собраны точные значения $d$ для проверенных
конструкций; каждая строка означает оценку $\chi(\RR^n;[1,\ell])\le k$ при
$\ell\le d$ (теорема~\ref{thm:main}).

\begin{table}[h!]
\centering
\caption{Проверенные интервальные константы:
$\chi(\RR^n;[1,\ell])\le k$ при $\ell\le d$.}
\label{tab:constants}
\medskip
\begin{tabular}{ccccc}
\hline
$n$ & решётка & $k$ & $d$ (точно) & $d$ (числ.)\\
\hline
$2$ & $A_2$            & $7$    & $\sqrt{7}/2$   & $1.322876$\\
$3$ & BCC              & $15$   & $1$            & $1.000000$\\
$3$ & FCC              & $21$   & $\sqrt{7/6}$   & $1.080123$\\
$4$ & $D_4$            & $49$   & $\sqrt{7/6}$   & $1.080123$\\
$4$ & $A_4^{*}$        & $54$   & $\sqrt{23/20}$ & $1.072381$\\
$5$ & $A_5^{*}$        & $140$  & $\sqrt{39/35}$ & $1.055597$\\
$6$ & $(3+\omega)E_6^{*}$ & $343$  & $\sqrt{7/6}$   & $1.080123$\\
$8$ & $(3+\omega)E_8$     & $2401$ & $\sqrt{7/6}$   & $1.080123$\\
\hline
\end{tabular}
\end{table}

Комментарии к таблице.

\begin{itemize}
\item Строки $n=2,3$ воспроизводят классические конструкции (Исбелл; Кулсон
      \cite{Coulson2002}; таблицы Иванова \cite{Ivanov2011}); строка BCC ---
      граничный случай $d=1$, реализуемый полуоткрытыми ячейками.
\item Для $D_4$ полный перебор всех подрешёток индексов $k=2,\dots,60$ (на
      индексе $49$ --- все $140\,050$ подрешёток) показал, что $49$ ---
      \emph{минимальное} допустимое $k$ для $D_4$ (и единственное в этом
      диапазоне); это согласуется с \cite[Prop.~12]{Arman2021}.
\item Для $A_5^{*}$ значение $d=\sqrt{39/35}$ оптимально среди \emph{всех}
      $1\,424\,115\,231$ подрешёток индекса $140$ (исчерпывающий
      HNF-перебор); исходная конструкция --- явная матрица $C_5$ из
      \cite{Arman2021}.
\item Точные значения $d$ для конструкций из \cite{Arman2021}
      (строки $k=140$, $343$, $2401$) ранее не публиковались; для
      $(3+\omega)$-конструкций доказанная в \cite{Arman2021} гарантия
      составляла лишь $d\ge 3/(2\sqrt{2})\approx 1.0607$, тогда как точное
      значение равно $\sqrt{7/6}=1.080123\ldots$
\end{itemize}

\subsection{Закон конструкции Эйзенштейна}

Пусть решётка $\Lambda$ несёт эйзенштейнову структуру (действие кольца
$\ZZ[\omega]$, $\omega=e^{2\pi i/3}$), и $\alpha=3+\omega$, $|\alpha|^2=7$.
Подрешётка $\alpha\Lambda$ имеет индекс $7^{n/2}$. Вычислительно установлено
(точно, в размерностях $n=2,4,6,8$):
\[
D(\alpha\Lambda)^2=\frac{7}{3}\,\lambda_1(\Lambda)^2,
\qquad\text{т.е.}\qquad
d=\frac{\sqrt{7/3}}{\,2R/\lambda_1\,},
\]
где $\lambda_1$ --- длина кратчайшего вектора $\Lambda$, а $R$ --- радиус
покрытия ($\diam V_0=2R$, замечание~\ref{rem:diam}): нормированный интервал
зависит только от отношения «покрытие/упаковка» решётки.

\begin{table}[h!]
\centering
\caption{Закон конструкции Эйзенштейна:
$d=\sqrt{7/3}\,/\,(2R/\lambda_1)$, $k=7^{n/2}$;
проверено точно в $n=2,4,6,8$.}
\label{tab:law}
\medskip
\begin{tabular}{ccccc}
\hline
$n$ & $\Lambda$ & $k=7^{n/2}$ & $2R/\lambda_1$ & $d$\\
\hline
$2$ & $A_2$     & $7$    & $2/\sqrt{3}$ & $\sqrt{7}/2$\\
$4$ & $D_4$     & $49$   & $\sqrt{2}$   & $\sqrt{7/6}$\\
$6$ & $E_6^{*}$ & $343$  & $\sqrt{2}$   & $\sqrt{7/6}$\\
$8$ & $E_8$     & $2401$ & $\sqrt{2}$   & $\sqrt{7/6}$\\
\hline
\end{tabular}
\end{table}

Для решёток с $2R/\lambda_1=\sqrt{2}$ (таковы $D_4$, $E_6^{*}$, $E_8$ и
решётка Лича $\Lambda_{24}$) закон даёт ровно $\sqrt{7/6}$, для $A_2$
($2R/\lambda_1=2/\sqrt{3}$) --- ровно $\sqrt{7}/2$.

\begin{remark}[следствие-гипотеза]
Если закон верен и для решётки Лича (прямая проверка вычислительно тяжела),
то $\chi\bigl(\RR^{24};[1,\sqrt{7/6}]\bigr)\le 7^{12}$.
\end{remark}

\section{Заключение}

Метод периодических раскрасок по решётчатым разбиениям Вороного даёт
вычислимые и сертифицируемые верхние оценки интервальных хроматических чисел:
полный HNF-перебор подрешёток, доказуемая граница перечисления кандидатов
(следствие~\ref{cor:window}) и самопроверка ячейки Вороного превращают каждое
найденное значение $d$ в проверяемый сертификат. Естественные продолжения ---
фактор-редукция перебора по группе автоморфизмов решётки и систематический
поиск в размерности $6$, где вопрос о минимальном числе цветов решёточных
раскрасок остаётся открытым \cite{Arman2021}.

\begin{thebibliography}{99}

\bibitem{Voronoi1952}
Вороной~Г.~Ф. Собрание сочинений в 3-х томах. --- Киев: Изд-во АН УССР, 1952.

\bibitem{Ivanov2006}
Иванов~Л.~Л. Об одной оценке хроматического числа пространства $\RR^4$ //
Успехи матем. наук. --- 2006. --- Т.~61, вып.~5. --- С.~181--182.

\bibitem{Ivanov2011}
Иванов~Л.~Л. О хроматических числах $\RR^2$ и $\RR^3$ с интервалами
запрещённых расстояний // Вестник РУДН, сер. Математика. Информатика.
Физика. --- 2011. --- №~1. --- С.~14--29.

\bibitem{Raigorodskii2000}
Райгородский~А.~М. О хроматическом числе пространства // Успехи матем.
наук. --- 2000. --- Т.~55, вып.~2. --- С.~147--148.

\bibitem{Raigorodskii2003}
Райгородский~А.~М. Хроматические числа. --- М.: МЦНМО, 2003.

\bibitem{Soifer2004}
Сойфер~А. Хроматическое число плоскости: его прошлое, настоящее и будущее //
Матем. просвещение. --- 2004. --- Вып.~8.

\bibitem{Arman2021}
Arman~A., Bondarenko~A., Prymak~A., Radchenko~D. Upper bounds on chromatic
number of $\mathbb{E}^n$ in low dimensions. --- 2021. ---
arXiv:2112.13438 [math.CO].

\bibitem{Borsuk1933}
Borsuk~K. Drei S\"atze \"uber die $n$-dimensionale euklidische Sph\"are //
Fundamenta Math. --- 1933. --- Vol.~20. --- P.~177--190.

\bibitem{Coulson2002}
Coulson~D. A 15-colouring of 3-space omitting distance one // Discrete
Math. --- 2002. --- Vol.~256. --- P.~83--90.

\bibitem{deGrey2018}
de~Grey~A.~D.~N.~J. The chromatic number of the plane is at least~5. ---
2018. --- arXiv:1804.02385 [math.CO].

\bibitem{DutourSikiric2021}
Dutour Sikiri\'c~M., Madore~D.~A., Moustrou~P., Vallentin~F. Coloring the
Voronoi tessellation of lattices // J. London Math. Soc. --- 2021. ---
Vol.~104, no.~3. --- P.~1135--1155.

\bibitem{Exoo2014}
Exoo~G., Ismailescu~D., Lim~M. On the chromatic number of $\RR^4$ //
Discrete Comput. Geom. --- 2014. --- Vol.~52, no.~2. --- P.~416--423.

\bibitem{Frankl1981}
Frankl~P., Wilson~R.~M. Intersection theorems with geometric consequences //
Combinatorica. --- 1981. --- Vol.~1. --- P.~357--368.

\bibitem{Hadwiger1944}
Hadwiger~H. Ein \"Uberdeckungssatz f\"ur den Euklidischen Raum //
Portugal. Math. --- 1944. --- Vol.~4. --- P.~140--144.

\bibitem{Jensen1995}
Jensen~T.~R., Toft~B. Graph Coloring Problems. --- New York:
Wiley-Interscience, 1995. --- P.~150--152.

\bibitem{Larman1972}
Larman~D.~G., Rogers~C.~A. The realization of distances within sets in
Euclidean space // Mathematika. --- 1972. --- Vol.~19. --- P.~1--24.

\bibitem{Nechushtan2002}
Nechushtan~O. On the space chromatic number // Discrete Math. --- 2002. ---
Vol.~256. --- P.~499--507.

\bibitem{Radoicic2003}
Radoi\v{c}i\'c~R., T\'oth~G. Note on the chromatic number of the space //
Discrete and Computational Geometry. The Goodman---Pollack Festschrift.
Algorithms and Combinatorics, vol.~25. --- Berlin: Springer, 2003. ---
P.~695--698.

\bibitem{Soifer1991}
Soifer~A. Chromatic number of the plane: a historical essay //
Geombinatorics. --- 1991. --- Vol.~1. --- P.~13--15.

\bibitem{Soifer2008}
Soifer~A. The Mathematical Coloring Book: Mathematics of Coloring and the
Colorful Life of Its Creators. --- New York: Springer, 2009. --- 607~p.

\end{thebibliography}

\end{document}
